\chapter{Übergangsfamilien und Klassifikation}

\section*{Motivation}
Übergangsbögen sind ein zentraler Bestandteil des Alignment-Entwurfs; ihre
Behandlung in der Literatur ist jedoch auf geometrische, ingenieurtechnische
und dynamische Sichtweisen verteilt. Für systematischen Vergleich, Bewertung
und Optimierung ist deshalb eine einheitliche Klassifikation erforderlich.

\section{Krümmungsbasierte Darstellung}
\begin{proposition}[Krümmungsdarstellung von Übergangsbögen]
\label{prop:transition-curvature-representation}
Jede reguläre ebene Kurve der Klasse \(C^2\), die auf \([0,L]\) nach der
Bogenlänge parametrisiert ist, bestimmt eine stetige Krümmungsfunktion
\[
\kappa:[0,L]\to\R ;
\]
umgekehrt bestimmt \(\kappa\) zusammen mit der Anfangspose die Kurve
eindeutig (\cref{prop:curvature-driven-reconstruction}). Übergangsbögen,
die diese Voraussetzungen erfüllen, werden daher bis auf die Anfangspose
durch ihre Krümmungsfunktion dargestellt.
\end{proposition}
\begin{proof}
Für eine \(C^2\)-Kurve mit Einheitsgeschwindigkeit gilt
\(\gamma'=(\cos\theta,\sin\theta)\) mit einer \(C^1\)-Richtung \(\theta\),
und \(\kappa=\theta'\) ist stetig \cite{DoCarmo1976Curves}. Die Umkehrung
ist \cref{prop:curvature-driven-reconstruction}.
\end{proof}
Die Voraussetzungen sind wesentlich: Eine Polylinie besitzt in ihren
Knoten keine klassische Krümmung, und eine beliebig parametrisierte Kurve
muss zunächst nach der Bogenlänge umparametrisiert werden. Unter diesen
Voraussetzungen ist die Darstellung von der ursprünglichen Kurvendefinition
unabhängig und stellt einen gemeinsamen mathematischen Rahmen bereit.

\subsection{Normierung}
\begin{definition}
Ein normierter Übergang ist auf
\[
s\in[0,1]
\]
definiert mit
\[
\kappa(0)=0,\qquad\kappa(1)=1.
\]
\end{definition}
Normierung entfernt den Maßstab und ermöglicht den unmittelbaren Vergleich von
Übergangsfamilien.

\section{Klassifikation nach funktionaler Form}
\subsection{Lineare Krümmung: Klothoide}
\begin{definition}
\[
\kappa(s)=as.
\]
\end{definition}
Die Klothoide ist durch die konstante Krümmungsableitung
\[
\kappa'(s)=\text{konstant}
\]
gekennzeichnet. Diese Einfachheit erklärt ihre weite Verbreitung im
Eisenbahningenieurwesen.

\subsection{Polynomiale Krümmungsfamilien}
\begin{definition}
\[
\kappa(s)=\sum_{i=0}^{n}a_i s^i.
\]
\end{definition}
Polynomiale Übergänge erlauben die ausdrückliche Regelung höherer
Krümmungsableitungen. Diese Freiheit ermöglicht Optimierung nach dynamischen
Kriterien. Übergänge des Klauder-Typs gehören zu dieser Klasse und beziehen
dynamische Erwägungen unmittelbar in den Krümmungsentwurf ein.

\subsection{Geometrische Familien}
Manche Übergangsbögen werden implizit oder durch geometrische Beziehungen
statt durch ausdrückliche Krümmungsfunktionen definiert. Zur Integration in
den vorliegenden Rahmen müssen sie in eine Krümmungsdarstellung überführt
werden.

\subsection{Engineering-basierte Familien}
Engineering-Ansätze definieren Übergangsbögen häufig durch Entwurfsregeln oder
empirische Erwägungen. Der Bloss-Übergang ist ein wichtiges Beispiel eines
solchen praktischen Engineering-Kompromisses.

\subsection{Optimierte und hybride Familien}
Moderne Ansätze führen Übergangsbögen ein, die ausdrücklich auf dynamische
Leistung optimiert sind. Dazu gehören Übergänge mit geglätteter Krümmung und
ruckoptimierte Polynomkurven.

\section{Klassifikation nach Verhalten höherer Ordnung}
\subsection{Krümmungsableitungen}
\begin{definition}
\[
\kappa'(s),\quad\kappa''(s)
\]
bezeichnen die erste und zweite Ableitung der Krümmung.
\end{definition}
Diese Größen bestimmen das dynamische Verhalten einschließlich
Beschleunigungsänderung und Glattheit der Kraftentwicklung.

\subsection{Glattheitsklassen}
\begin{definition}
Ein Übergang gehört zur Klasse \(C^k\), wenn \(\kappa\) \(k\)-mal stetig
differenzierbar ist.
\end{definition}
Höhere Glattheit verbessert typischerweise die dynamische Leistung und
verringert den Verschleiß.

\section{Vergleich der Familien}
\begin{center}
\begin{tabular}{lccc}
\textbf{Familie} & $\kappa$ & $\kappa'$ & $\kappa''$ \\
\hline
Klothoide & linear & konstant & 0 \\
Polynomial & flexibel & flexibel & flexibel \\
Engineering & implizit & implizit & implizit \\
Optimiert & entworfen & kontrolliert & kontrolliert \\
\end{tabular}
\end{center}

\section{Interpretation als Steuerungsfunktionen}
\begin{proposition}
Übergangsbögen können als Steuerungsfunktionen interpretiert werden, die durch
Krümmungsentwicklung auf die Systemdynamik wirken.
\end{proposition}
Diese Interpretation verbindet geometrischen Entwurf unmittelbar mit
dynamischer Leistung.

\section{Verbindung zum Eisenbahningenieurwesen}
Der Entwurf von Eisenbahnübergängen wird durch dynamische Bedingungen wie
Beschleunigung, Ruck und Verschleiß bestimmt. Diese Bedingungen werden
unmittelbar zu Anforderungen an
\[
\kappa(s),\quad\kappa'(s),\quad\kappa''(s).
\]

\section{Zentrale Erkenntnis}
\begin{theorem}
Übergangsbögen sind nicht lediglich geometrische Verbindungen, sondern
Steuerungsfunktionen, die das Systemverhalten bestimmen.
\end{theorem}

\section{Schlussfolgerung}
\begin{summary}
Übergangsbögen bilden eine einheitliche Klasse von Krümmungsfunktionen.

Ihre Klassifikation nach funktionaler Form und Verhalten höherer Ordnung
stellt einen systematischen Rahmen für Vergleich, Optimierung und Integration
in Engineering-Systeme bereit.

Diese Sicht schafft eine unmittelbare Verbindung zwischen Geometrie und
Systemdynamik.
\end{summary}
